\section{Evaluation}
\label{sec:evaluation}

We split the research question into three sub-questions.
\begin{description}
  \item[RQ1 (soundness)] Does the checker ever give a wrong verdict in practice?
  \item[RQ2 (coverage)] How much of the input domain does it certify, compared with one SMT
        query over the whole domain and with differential testing?
  \item[RQ3 (layers)] What does each layer contribute? What do the intervals add on top of
        the solver, and the solver on top of the intervals, and at what cost?
\end{description}

\subsection{Setup}
\label{sec:setup}

\paragraph{Benchmarks}
We use 20 method pairs in three groups. The \emph{corpus} (\code{Refactorings.java}) holds 10
pairs that we wrote for this project. Five are preserving edits: inlining a temporary,
removing dead code, inverting a branch, reversing a \code{max} comparison, and an unchanged
pair. Three are value-only faults: \code{sign}, \code{absVal} and an off-by-one comparison
\code{belowTwo}. Two are exception faults: guard removals in \code{ratio} and
\code{reciprocal}. For the corpus, ground truth comes from the real JVM. A Java harness runs
both versions with \code{-ea} on all 238 inputs of the declared domains and records outcome
and return value. Five pairs
come from the stock \jpamb{} class \code{Simple}. Two of them compare a method with itself and
act as controls. Five evaluation-only pairs (\code{EvalPairs.java}) contain loops or many
paths. The loop pairs are \code{loopDiv} with an added guard, \code{loopDiv} with a rewritten
body, and \code{loopThenZero}. The many-path pairs are two versions of \code{ladder}, which has
twelve sequential \code{if} statements. One version reorders them and the other changes one
constant.

\paragraph{Configurations}
We compare four configurations on each pair.
\begin{itemize}
  \item \code{gate}: intervals, verdict rules 0 to 2b$'$ and refinement, without the solver.
  \item \code{hybrid}: our full approach.
  \item \code{z3only}: one SMT query over the whole domain. It encodes both the error paths
        and the returning paths, so it compares outcome class and value together. This is a
        baseline in the style of differential symbolic execution~\citep{person2008dse},
        without the interval layer.
  \item \code{random}: differential testing~\citep{mckeeman1998differential} with 1000
        uniformly random inputs (seed 0) on a concrete interpreter with Java \code{int}
        semantics. Without a difference it answers \unk{} and never \pres{}.
\end{itemize}
We run every pair on a \emph{small} domain $[-3,3]^k$ and on a \emph{large} domain
$[-1000,1000]^k$, which has up to 4 million inputs. On the small domain ground truth comes
from exhaustive execution. On the large domain we check a \pres{} verdict against 2000 random
inputs. We replay solver witnesses (Rule~3$'$). We do not replay whole-region \chg{} verdicts.

\paragraph{Metrics}
A verdict is \emph{wrong} if it says \pres{} while a difference exists, or \chg{} without a
real difference. An \unk{} verdict is never wrong. A pair is \emph{decided} if the verdict is
not \unk{}. The \emph{certified coverage} is the fraction of the domain with a proved verdict
(\cref{sec:formal}). For the baselines it is $1$ for \pres{} and one input for a witness.

\paragraph{Environment}
We use Python~3.13, the \jpamb{} 0.9.0 bytecode model and Z3 5.4.1, with a 5\,s timeout per
SMT query and a region budget of about 200. We run on a laptop (Intel Core i5-1038NG7, 16\,GB,
macOS). The times come from a single run. \Cref{app:variance} has repeated timings.

\subsection{Results}
\label{sec:results}

\begin{table}
  \centering
  \caption{Summary over 20 pairs per domain (40 runs per configuration). \emph{Coverage}
  is the mean certified coverage. \emph{SMT} is the number of solver queries in total.}
  \label{tab:results}
  \small
  \begin{tabular}{@{}l rrrr rrrr r@{}} \toprule
    & \multicolumn{4}{c}{small $[-3,3]^k$} & \multicolumn{4}{c}{large $[-1000,1000]^k$} & all \\
    \cmidrule(lr){2-5} \cmidrule(lr){6-9}
    Config & decided & coverage & SMT & ms & decided & coverage & SMT & ms & wrong \\ \midrule
    gate   & 19/20 & 98\% & 0   & 1014 & 13/20 & 61\% & 0   & 4418  & 0 \\
    hybrid & 19/20 & 98\% & 101 & 1120 & \textbf{19/20} & \textbf{95\%} & 566 & 6396 & 0 \\
    z3only & 17/20 & 49\% & 20  & 1082 & 17/20 & 40\% & 20  & 1265  & 0 \\
    random & 9/20  & 5\%  & 0   & 1165 & 9/20  & $<$1\% & 0 & 35422 & 0 \\ \bottomrule
  \end{tabular}
\end{table}

\Cref{tab:results} summarises the 160 runs. \Cref{tab:perpair} in \cref{app:perpair} lists
every pair.

\paragraph{RQ1: Soundness}
No configuration produced a wrong verdict. We checked each pair against exhaustive execution
on the small domain and against 2000 samples on the large one. Solver witnesses replay by
construction, because we report a model that does not replay as \unk{}. The pair harness, with a
matched 5\,s budget per pair, passes all 10 corpus pairs for both the analyser and seeded
differential testing. The pair \code{ladder/\_changed} shows that a verdict holds only for its
domain. The edit changes a constant that lies outside $[-3,3]$. The pair is therefore \pres{}
on the small domain and \chg{} on the large one. All four configurations agree with that.

\paragraph{RQ2: Coverage}
On the large domain the hybrid decides 19 of 20 pairs and certifies 95\% of the inputs on
average. One SMT query over the domain decides 17 of 20 and certifies 40\%. Random testing
decides 9 of 20 and certifies under 1\%. The gap in coverage comes from the \chg{} pairs. The
\code{z3only} and \code{random} configurations certify only their witness. The hybrid also
proves the rest of the domain. For \code{ratio} it proves 48 of 49 inputs \pres{} on the small
domain. On the large domain it proves all but the one input $(0,0)$ of four million. Random
testing finds this input in none of 20 seeds (\cref{app:variance}).

\paragraph{RQ3: Loops}
The symbolic executor rejects backward jumps, so \code{z3only} answers \unk{} on all three
loop pairs. The interval layer still decides two of them. Both versions of
\code{loopThenZero} end in \code{s / 0} after a counting loop that provably terminates. Their
outcome set is $\set{\dbz}$, and Rule~1 proves the pair \pres{} on the whole domain in one
region analysis. For \code{loopDiv} with an added guard, we split at $d = 0$. This isolates a
region with disjoint outcome sets, which gives a whole-region \chg{} verdict. It certifies
76\% (small) and 50\% (large) of the inputs.

\paragraph{RQ3: Cost}
On loop-free pairs one SMT query is the fastest way to a verdict. The \code{z3only}
configuration needs 1.3\,s for all large-domain runs, and the hybrid needs 6.4\,s. The extra
time buys coverage. After the hybrid finds a witness, it keeps splitting and querying. This
costs most in time on \code{justAdd/justMulitply} (144 queries, 2.7\,s) and on
\code{ladder/\_changed}. The two \code{loopDiv} pairs issue 190 queries each that fail at once
on loops.

\subsection{Discussion}
\label{sec:discussion}

On these 20 small, \code{int}-only pairs the answer to the research question is yes, with a
qualification. No verdict was wrong on any run. The hybrid is the only configuration that
both decides most pairs and says where an edit changes behaviour. The two layers complement
each other. The intervals decide regions where the outcome classes differ or coincide on a
single error, and they handle loops. The solver decides regions where both versions return and
only the value is in question. Alone, each layer falls short. Without the solver, the gate
decides none of the five preserving corpus pairs on the large domain, because intervals cannot
relate the two versions. Without the gate, \code{z3only} cannot handle loops and certifies no
more than a witness. The qualification is that \code{z3only} is the faster way to a yes or no
verdict on loop-free code. The many-path pairs caused no path explosion, so we make no claim
about scalability in the number of paths. The pair \code{loopDiv/\_rewritten} stays \unk{}
everywhere. Both versions return a loop-dependent value that neither layer can relate.

\subsection{Threats to Validity}
\label{sec:threats}

\paragraph{Internal}
Coverage favours methods that localise, because \code{z3only} stops at one witness. We
therefore treat decided runs as the primary measure. We wrote the corpus and the extra pairs
ourselves. The splitting heuristic uses the constants that appear in the code, which suits
small hand-written pairs. We reduce the risk of bugs in the checker in three ways. We take the
corpus ground truth from the real JVM, we replay solver witnesses, and we test the interval
operations against Java semantics with property-based tests. For the \code{Simple} and
\code{EvalPairs} groups the ground truth comes from our Python interpreter. We tested that
interpreter against the real-JVM observations, but only on the corpus. On the large domain we
only spot-check a \pres{} verdict by sampling, so a soundness bug that affects few inputs could
go unnoticed there. We did not mechanise the proofs in \cref{sec:guarantees}. Code review found
one gap between proof and code. Exception tables were ignored, so a caught division gave a
false \chg{} verdict. We now reject methods with handlers. Soundness assumes \code{-ea} and
\code{javac}-shaped loops. A 10,000-step cap gives \emph{unsupported}.

\paragraph{External}
The 20 pairs are small, static, \code{int}-only and call-free. The domains are boxes of at
most $2001^2$ inputs. Real refactorings involve objects, calls and caught exceptions, which our
observation contract excludes. The results say nothing about them. A verdict holds only on the
declared domain, not on all of \code{int}.
